\section{CONCLUDING REMARKS AND FUTURE DIRECTIONS}  \label{sec:con}
In this work, we introduce \qs, the first theoretically principled
approach to pruning large ground sets for constrained discrete maximization
problems. Moreover, the theoretical guarantees extend to a range of budgets so
that the pruned ground set is applicable to more than a single optimization problem.
In addition to the theoretical properties, our algorithm exhibits excellent empirical
performance, and outperforms pre-existing heuristics for the pruning problem.
Future directions include improving the theoretical properties, as our ratio $\alpha$,
while constant, is rather small. Also, upper bounds on the pruning problem would shed
light on the complexity of the problem. Currently, we are unaware of existing hardness
results that would imply that even a ratio of $\alpha = 1$ is impossible.
Finally, extending the framework to handle more than one objective function, as well
as deletions to the ground set, would be an interesting avenue for future work. 
\clearpage

%%% Local Variables:
%%% mode: latex
%%% TeX-master: "main.tex"
%%% End:
